\documentclass[../../main.tex]{subfiles}
\graphicspath{{../../image/}{../../figure/}}
\begin{document}

\section{混合密度网络}\label{section:ch06-混合密度网络}

到目前为止，我们已经讨论了输出表示简单概率分布的神经网络，包括连续变量的高斯分布以及离散变量的二元分布。在本章最后一节，我们将展示神经网络如何通过将网络输出视为更复杂分布（如高斯混合分布）的参数来表示更一般的条件概率。通过这种方式定义关联的误差函数和相应输出单元激活函数，我们称此类网络为\textbf{混合密度网络}（mixture density network）。

\subsection{机器人运动学示例}\label{section:ch06-机器人运动学示例}

监督学习的目标是建模条件分布 $p(t\mid \mathbf{x})$。对于许多简单回归问题，选择高斯分布较为常见；然而实际机器学习问题往往具有明显非高斯分布特征。例如在分布可能是多模态的\textbf{逆问题}（inverse problem）中，高斯假设可能导致非常糟糕的预测。

作为逆问题的简单说明，考虑机械臂的运动学（图~6.16）。\textbf{正问题}（forward problem）涉及在给定关节角度时找到末端执行器位置，并具有唯一解；然而在实践中，我们希望将末端执行器移动到特定位置，为此必须设置适当的关节角度。我们需要求解一个逆问题，如图~6.16~所示，该问题有两个解（见习题~6.16），对应于"肘部向上"和"肘部向下"两种配置。

正问题通常对应于物理系统中的因果关系，通常具有唯一解——例如人体中特定症状可能由特定疾病引起。然而在机器学习中，我们通常需要求解逆问题，例如在给定一组症状情况下预测疾病。如果正问题涉及多对一映射，则逆问题将具有许多解：例如多种不同疾病可能导致相同症状。

在机器人运动学示例中，运动学由几何方程定义，其多模态性显而易见。但在许多机器学习问题中，多模态的存在可能不太明显，特别是涉及高维空间的问题。出于教学目的，下面考虑一个简单小问题，其多模态可以很容易可视化。该问题数据通过在区间 $(0,1)$ 上均匀采样变量 $x$ 产生，得到一组值 $\{x_n\}$；通过计算函数 $x+0.3\sin(2\pi x)$，然后添加在区间 $(-0.1,0.1)$ 上均匀采样得到的噪声，得到相应目标值。

图~6.17~左图展示了正问题数据集，红色曲线显示了通过最小化平方和误差函数拟合两层神经网络的结果（该网络具有~6~个隐藏单元和一个线性输出单元）。最小二乘对应于高斯假设下的最大似然：它对正问题是一个很好的模型，但对高度非高斯的逆问题却是一个很差的模型。图~6.17~右图显示了通过交换 $x$ 和 $t$ 角色得到的逆问题——再次通过最小化平方和误差函数训练同样的网络，由于数据的多模态性，网络对数据的拟合效果较差。

\subsection{条件混合分布}\label{section:ch06-条件混合分布}

我们想要找到建模条件概率分布的一般框架。这可以通过使用专为建模 $p(t\mid \mathbf{x})$ 设计的混合模型来实现，其中混合系数和分量密度都是输入向量 $\mathbf{x}$ 的灵活函数，从而得到\textbf{混合密度网络}（mixture density network）。对于任何给定 $\mathbf{x}$ 值，混合模型为建模任意条件密度函数 $p(t\mid \mathbf{x})$ 提供一般形式；如果有足够灵活的网络，我们就有了一个可以近似任意条件概率分布的框架。

在这里我们显式地为高斯分量设计模型如下：
\begin{align}
p(t\mid \mathbf{x}) = \sum_{k=1}^{K}\pi_k(\mathbf{x})\,\mathcal{N}\!\left(t\,\middle|\,\mu_k(\mathbf{x}),\sigma_k^{2}(\mathbf{x})\right).
\label{align:ch06-mdn}
\end{align}

除了高斯分布，还可以为分量使用其他分布：例如如果目标变量是二元变量而非连续变量，则可以使用伯努利分布。虽然这里使用各向同性协方差，但也可以通过使用 Cholesky 因式分解表示协方差来使混合密度网络支持一般协方差矩阵~\citep{Williams1996}。即使使用各向同性分量，由于混合分布的影响，条件分布 $p(t\mid \mathbf{x})$ 也不具备对 $t$ 各分量的可分解性（与标准平方和回归模型不同）。

令混合模型的各个参数——即混合系数 $\pi_k(\mathbf{x})$、均值 $\mu_k(\mathbf{x})$ 和方差 $\sigma_k^{2}(\mathbf{x})$——由以 $\mathbf{x}$ 为输入的神经网络的输出决定。混合密度网络结构如图~6.18~所示。混合密度网络与混合专家模型密切相关~\citep{Jacobs1991}：主要区别在于混合专家模型对自身每个分量模型都有独立参数，而在混合密度网络中，我们使用相同的函数来预测所有分量密度的参数和混合系数，于是非线性隐藏单元在依赖于输入的函数之间是共享的。

如果混合模型~\eqref{align:ch06-mdn}~中有 $K$ 个分量，并且 $t$ 有 $L$ 个分量，则网络将具有：$K$ 个输出单元预激活（用 $a_k^{\pi}$ 表示）用于确定混合系数 $\pi_k(\mathbf{x})$；$K$ 个输出（用 $a_k^{\sigma}$ 表示）用于确定高斯标准差 $\sigma_k(\mathbf{x})$；$K\times L$ 个用 $a_{kj}^{\mu}$ 表示的输出用于确定高斯均值分量 $\mu_{k,j}(\mathbf{x})$。网络输出总数由 $L+2K$ 给出，这与具有 $L$ 个输出、用于预测目标变量条件均值的网络不同。例如图~6.18~中的神经网络可以是具有 S 形 tanh 隐藏单元的两层网络。

混合系数必须满足下列约束：
\begin{align}
\sum_{k=1}^{K}\pi_k(\mathbf{x})=1, \qquad 0\le \pi_k(\mathbf{x})\le 1,
\label{align:ch06-pi-constraint}
\end{align}
这可以使用一组 softmax 输出来实现：
\begin{align}
\pi_k(\mathbf{x}) = \frac{\exp(a_k^{\pi})}{\sum_{j=1}^{K}\exp(a_j^{\pi})}.
\label{align:ch06-pi-softmax}
\end{align}

同样，方差必须满足 $\sigma_k^{2}(\mathbf{x})\ge 0$，因而 $\sigma_k(\mathbf{x})$ 可以使用相应网络预激活的指数来表示：
\begin{align}
\sigma_k(\mathbf{x}) = \exp(a_k^{\sigma}).
\label{align:ch06-sigma-exp}
\end{align}

最后由于均值 $\mu_k(\mathbf{x})$ 具有实分量，可以用网络输出直接表示：
\begin{align}
\mu_{k,j}(\mathbf{x}) = a_{kj}^{\mu},
\label{align:ch06-mu-linear}
\end{align}
其中输出单元激活函数由恒等式 $f(a)=a$ 给出。

\subsection{梯度优化}\label{section:ch06-梯度优化}

混合密度网络的可学习参数包括神经网络中的权重和偏置向量 $\mathbf{w}$，它可以通过最大似然来确定，也可以等价地通过最小化定义为似然负对数的误差函数来确定。对于彼此独立的数据，误差函数采用以下形式（其中显式给出了关于 $\mathbf{w}$ 的依赖）：
\begin{align}
E(\mathbf{w}) = -\sum_{n=1}^{N}\ln\!\left\{\sum_{k=1}^{K}\pi_k(\mathbf{x}_n,\mathbf{w})\,\mathcal{N}\!\left(t_n\,\middle|\,\mu_k(\mathbf{x}_n,\mathbf{w}),\sigma_k^{2}(\mathbf{x}_n,\mathbf{w})\right)\right\}.
\label{align:ch06-mdn-nll}
\end{align}

为了最小化误差函数，需要计算误差 $E(\mathbf{w})$ 关于参数 $\mathbf{w}$ 的导数。稍后我们将看到如何自动计算这些导数（参见第~8~章）。然而显式地为输出单元预激活的误差导数导出合适表达式是有帮助的，因为这强调了这些量的概率解释。由于误差函数~\eqref{align:ch06-mdn-nll}~由一些项的和组成，其中每项对应一个训练数据点，因此可以考虑特定输入向量 $\mathbf{x}_n$ 与相关目标向量 $t_n$ 的导数；总误差 $E$ 的导数通过对所有数据点求和获得，每个数据点自身的梯度也可以直接用于基于梯度的优化算法（参见第~7~章）。

引入以下变量将有助于推导：
\begin{align}
\gamma_{nk} = \frac{\pi_k(\mathbf{x}_n)\,\mathcal{N}\!\left(t_n\,\middle|\,\mu_k(\mathbf{x}_n),\sigma_k^{2}(\mathbf{x}_n)\right)}{\sum_{j=1}^{K}\pi_j(\mathbf{x}_n)\,\mathcal{N}\!\left(t_n\,\middle|\,\mu_j(\mathbf{x}_n),\sigma_j^{2}(\mathbf{x}_n)\right)},
\label{align:ch06-gamma}
\end{align}
其中 $\mathcal{N}_k$ 表示 $\mathcal{N}(t_n\mid \mu_k(\mathbf{x}_n),\sigma_k^{2}(\mathbf{x}_n))$。这些量可被自然解释为混合分布分量的后验概率，其中混合系数 $\pi_k(\mathbf{x})$ 可以视为与 $\mathbf{x}$ 相关的先验概率（见习题~6.17）。

关于控制混合系数的网络输出预激活 $a_k^{\pi}$ 的误差函数导数由下式给出（参见习题~6.18）：
\begin{align}
\frac{\partial E_n}{\partial a_k^{\pi}} = \pi_k(\mathbf{x}_n) - \gamma_{nk}.
\label{align:ch06-mdn-deriv-pi}
\end{align}

类似地，关于控制分量均值的网络输出预激活 $a_{kj}^{\mu}$ 的误差函数导数由下式给出（见习题~6.19）：
\begin{align}
\frac{\partial E_n}{\partial a_{kj}^{\mu}} = \gamma_{nk}\frac{\mu_{k,j}(\mathbf{x}_n)-t_{n,j}}{\sigma_k^{2}(\mathbf{x}_n)}.
\label{align:ch06-mdn-deriv-mu}
\end{align}

最后，关于控制分量方差的网络输出预激活 $a_k^{\sigma}$ 的误差函数导数由下式给出（见习题~6.20）：
\begin{align}
\frac{\partial E_n}{\partial a_k^{\sigma}} = \gamma_{nk}\left\{\frac{\|\mathbf{t}_n-\boldsymbol{\mu}_k(\mathbf{x}_n)\|^2}{\sigma_k^{2}(\mathbf{x}_n)}-1\right\}.
\label{align:ch06-mdn-deriv-sigma}
\end{align}

\subsection{预测分布}\label{section:ch06-预测分布}

混合密度网络训练好后，给定任意输入向量值，就可以预测目标数据的条件密度函数。就预测输出向量值的问题而言，该条件密度表示对数据生成器的完整描述。从这个密度函数中，我们可以计算出不同应用可能关注的更具体量。其中最简单的一个量是平均值，它对应于目标数据的条件平均：
\begin{align}
\mathbb{E}[t\mid \mathbf{x}] = \int t\,p(t\mid \mathbf{x})\,dt = \sum_{k=1}^{K}\pi_k(\mathbf{x})\,\mu_k(\mathbf{x}).
\label{align:ch06-mdn-condmean}
\end{align}

由于使用最小二乘法训练的标准网络会近似条件均值，因此我们看到混合密度网络可以作为一个特例再现传统最小二乘结果。当然正如我们已经指出的，对于多峰分布，条件均值的值是有限的。

我们可以类似地评估关于条件均值密度函数的方差（见习题~6.21）：
\begin{align}
\mathbb{V}[t\mid \mathbf{x}] = \mathbb{E}[\|t-\mathbb{E}[t\mid \mathbf{x}]\|^2\mid \mathbf{x}],
\label{align:ch06-mdn-condvar-def}
\end{align}
展开后为
\begin{align}
\mathbb{V}[t\mid \mathbf{x}] = \sum_{k=1}^{K}\pi_k(\mathbf{x})\left\{\sigma_k^{2}(\mathbf{x})+\|\boldsymbol{\mu}_k(\mathbf{x})-\mathbb{E}[t\mid \mathbf{x}]\|^2\right\}.
\label{align:ch06-mdn-condvar}
\end{align}

这里用到了~\eqref{align:ch06-mdn}~和~\eqref{align:ch06-mdn-condmean}。该结果比相应的最小二乘结果更普遍，因为方差是关于 $\mathbf{x}$ 的函数。这是一个\textbf{异方差}（heteroscedastic）模型的示例：数据上的噪声方差是输入向量 $\mathbf{x}$ 的函数。

让我们回到图~6.17~所示的逆问题示例来说明混合密度网络的使用方法。图~6.19~显示了混合系数 $\pi_k(\mathbf{x})$、均值 $\mu_k(\mathbf{x})$ 和对应于 $p(t\mid \mathbf{x})$ 的条件密度等值线。神经网络的输出（即混合模型中的参数）必然是输入变量的连续单值函数；然而从图~6.19(c)~可以看到，该模型能够通过调制混合分量 $\pi_k(\mathbf{x})$ 的振幅来产生条件密度，该密度对 $x$ 的某些值是单峰的，而对其他值是三峰的。该模型有~3~个高斯分量，并使用一个两层神经网络，该网络有~5~个 S 形 tanh 隐藏单元和~9~个输出（分别对应于~3~个高斯分量的均值、~3~个方差以及~3~个混合系数）。在 $x$ 取小值或大值时，目标数据条件概率密度是单峰的，只有一个高斯分量先验概率很高；而在 $x$ 取中间值时，条件概率密度是三峰的，~3~个混合系数具有差不多的值。

我们已经看到，对于多峰分布，条件均值给出了较差的数据表示。例如在图~6.16~所示简单机械臂控制中，我们需要从两种可能关节角度设置中选择一种来获得所需末端执行器位置，但两种解的平均值本身并不是一个解。在这种情况下，条件众数可能更有价值。由于混合密度网络的条件众数没有简单解析解，因此需要进行数值迭代。一种简单替代方法是在每个 $x$ 值处取最可能分量（即混合系数最大的分量）的平均值。这一做法已在图~6.19(d)~所示示例数据集上给出。

\end{document}
